% Transcription — fonds Alexandre Grothendieck, shelfmark 10, batch 10 (pages 181-197).
% Established from the facsimile archives/batches/10/batch-10.pdf (a mirror of
% https://grothendieck.umontpellier.fr/10.pdf), per the skill
% transcribe-grothendieck.
%
% Pass: Opus 5 (claude-opus-5), 2026-09-12 — first pass, unchecked against the
% pages by a human.
%
% The last seventeen leaves of the folder, all of them his own hand and all of
% them on the subject of the carton. Pages 181-182 finish without a break the
% run begun page 180 — the free group Γ = Z(I) on the indecomposable classes,
% ordered by Krull-Schmidt — and turn it into a ring: a tensor product on the
% class makes Γ a commutative algebra, the privileged object its unit, duality
% an automorphism, and the three together make I « un hypergroupe discret ».
% Then three worked cases, the last two numbered by him II and III; the first
% carries no number on the leaves.
%
% Pages 183-188, the additive group. The rational representations of k are
% unipotent, so their classification is that of nilpotent operators, i.e. of the
% X-primary k[X]-modules of finite type, and I ≃ N⁺ with k[X]/X^i answering i.
% Pages 184-185 prove there is only the one indecomposable of dimension 1 — a
% polynomial P with P(x+y) = P(x)P(y) has P(x+y) of degree n and P(x)P(y) of
% degree 2n, absurd — and set against it the characteristic-p count, where the
% f(λ,μ)(x) = λx^{p^h} + μ give infinitely many inequivalent classes E_h(y).
% Page 186 settles characteristic zero: Prop. 2, the classification of the
% rational representations of the abelian algebraic groups is that of the
% X-primary k[X]-modules of finite type. Pages 187-188 compute the ring: ξ = E_2
% the class of k[X]/(X²), ξ^n = Σ c_{ni} E_i with c_{nn} = 1, whence Prop. 3,
% A(k) = Z[ξ] the polynomial ring on one generator of degree 2 — the iterated
% T^k acting on F_n = k[X_1,…,X_n]/(X_1²,…,X_n²) as multiplication by
% Ṡ^k = k! Σ_{i_1<⋯<i_k} Ẋ_{i_1}⋯Ẋ_{i_k} is what carries it.
%
% Pages 189-190, « II  A(k*) ». The irreducible representations go through a
% rational P : k* → L*, P(λ) = Σ λ^i c_i, and P(λλ') = P(λ)P(λ') forces
% c_i = c_i², c_i c_j = 0, so a single c_n = 1 and P(λ) = λ^n: the only
% irreducibles are the σ_n of dimension 1, with σ_n σ_m = σ_{n+m} and
% σ̌_n = σ_{-n}. Every rational representation is then completely reducible and
% A(k*) ≃ Z(Z), which generalises to a torus as A(G) ≃ Z(Γ).
%
% Pages 191-197, « III  G = k · k* », the semi-direct product with k* acting by
% σ_p. In characteristic zero the Lie algebra is [Y,X] = pX; the weight
% decomposition V = Σ V_n forces X_{ji} = 0 but for j − i = p, indecomposability
% forces the non-zero V's to be the V_{n+pk}, and Lemmas 1 and 2 make the
% u_i = X_{n+p(i+1),n+pi} injective and then bijective, the E_i of dimension 1.
% The indecomposables are therefore the e_{n,k} (n ∈ Z, k ≥ 1), with
% e_{n,k} = σ_n ⊗ e_{0,k}, e_{0,k} the representation on k[Z]/Z^k, and
% ě_{n,k} = e_{−n−p(k−1),k}. Pages 195-197 compute A(G): Λ = A(H) = Z(Γ) sits
% inside it by j*(σ_n) = e_{n,1}, A(F) = Z[e_2] receives it by i*(e_{n,k}) = e_k,
% and a weight argument on ξ = e_{0,2} closes the folder on A(G) = Λ[ξ], the
% polynomial algebra over Λ on one generator with d(ξ) = 2 and ξ̌ = σ_{-p}·ξ.
%
% Three material facts govern the batch and are stated here rather than page by
% page. First, this run is drafting and not copying out: the formulas, the
% underlined statements and the displayed tables carry, and a great deal of the
% connective prose does not — pages 181-186 and 189-197 are written at speed,
% and the apparatus is dense in consequence. Pages 187-188 are the exception and
% read almost end to end. Second, three blocks are cancelled whole — the boxed
% middle of page 184 under two long diagonals, the head of page 193 under some
% fifteen strokes, and the two struck lines at the foot of page 196 — and they
% give up their formulas and little of the prose between them. Third, notes of
% his run up the left edge of pages 182, 183, 184, 186, 189 and 194 in a hand
% faster still; they are given as far as they go and no further.
%
%   181  the dimension function on Γ⁺, n(e) = 1 and n(ǐ) = n(i); what the data of
%        I, n and the Hom_{ji} amount to; the tensor product making Γ a
%        commutative algebra with e its unit and ∨ an automorphism; i * ǐ always
%        has a unit component along e, so I is « un hypergroupe discret ».
%   182  if every M of C is completely reducible and C comes from a group G, Γ
%        is the algebra of the linear combinations of the representations, and Γ⁺
%        the functions of the linear representations of G, all the structures
%        carried over (n(f) ↦ f(e), ∨ ↦ ∨, e ↦ 1); for G semi-simple, the
%        algebra of functions on T/W; the case G = k*, Γ(G) ≃ the group algebra
%        of Z with I(G) ≃ Z.
%   183  G = k: the rational representations are unipotent, hence the nilpotent
%        operators, hence the X-primary k[X]-modules; the elementary types X^k,
%        I ≃ N⁺ with k[X]/X^i, the unit 1, the involution the identity, « Quelle
%        est la multiplication ? »; marginal « Caractéristique 0 ? ».
%   184  the commutant of the P(x) is a field D, the P(x) take their values in
%        the centre Z, and P(x+y) = P(x)P(y) for a polynomial of degree n > 0
%        gives degree n against degree 2n, absurd; the boxed block cancelled.
%   185  so P(x) = 1, Z = k, and there is only the indecomposable of dim 1;
%        Corollaire: every rational representation is triangulable with 1's on
%        the diagonal; in characteristic p, f(λ,μ)(x) = λx^{p^h} + μ,
%        τ_y f(λ,μ) = f(λ, μ + λy^{p^h}), the representation E_h(y), and
%        infinitely many pairwise inequivalent classes against two in car. 0.
%   186  characteristic zero: the rational representations of the abelian groups
%        are the nilpotent representations of the abelian Lie algebra k;
%        Prop. 2, their classification is that of the (X)-primary k[X]-modules of
%        finite type; the indecomposables are the k[X]/X^n; Corollaire on A(k);
%        marginals « k de car⁰ » and the note on the characteristic.
%   187  E_n the class of k[X]/X^n, ∨ the identity, d(E_n) = n, E_1 the unit;
%        ξ = E_2, ξ^n the class of F_n = k[X_1,…,X_n]/(X_1²,…,X_n²) with T the
%        multiplication by Ṡ; T^k as multiplication by
%        Ṡ^k = k! Σ_{1≤i_1<⋯<i_k≤n} Ẋ_{i_1}⋯Ẋ_{i_k}; T^n F_n of dim 1 and
%        T^{n+1} = 0, so ξ^n is a Z⁺-combination of E_1,…,E_n with E_n's
%        coefficient 1.
%   188  ξ^n = L_n(E_1,…,E_n) = Σ c_{ni} E_i with c_{nn} = 1; Proposition 3,
%        A(k) = Z[ξ] with d(ξ) = 2; the table of the first five relations both
%        ways; the parenthetical on the hidden ring structure.
%   189  II, A(k*): the irreducibles go through P(λ) = Σ_{-n}^{n} λ^i c_i with
%        P(λλ') = P(λ)P(λ'), whence c_i = c_i², c_i c_j = 0, a single c_n = 1,
%        P(λ) = λ^n and L = k; the Proposition, the σ_n of dim 1 with
%        σ_n(λ)t = λ^n t; Corollaire σ_n σ_m = σ_{n+m}, σ̌_n = σ_{-n}.
%   190  Σ c_i = P(1) = 1, the c_i two-by-two orthogonal projectors; the
%        Proposition that every rational representation of k* is completely
%        reducible; Corollaire A(k*) ≃ Z(Z) with σ_n answering n; the
%        generalisation to tori, Γ ≃ Z^n and A(G) ≃ Z(Γ).
%   191  III, G = k · k* semi-direct, k* acting by σ_p(λ)t = λ^p t; the Lie
%        algebra extension of F = k by H = k with (1) [Y,X] = pX; V = Σ V_n with
%        Y·a = na, the components X_{ji} ∈ L(V_i,V_j) and (j−i)X_{ji} = pX_{ji},
%        so X_{ji} = 0 but for j − i = p.
%   192  if all the X_{ji} vanish, one V_n only, of dimension one, and the
%        representation is σ_n through G → H; otherwise V = V' + V'' with
%        V' = Σ_{n≡α(p)} V_n, and V' = V; the non-zero V's are the V_{n+pk},
%        k the largest with V_{n+pk} ≠ 0; u_i = X_{n+p(i+1),n+pi}, E_i = V_{n+pi},
%        the chain E_0 → E_1 → ⋯ → E_k, and Lemme 1, the u_i are injective.
%   193  the induction for Lemme 1: x in Ker u_ℓ, the supplements E'_j and E''_j,
%        E_j = E'_j ⊕ E''_j, ΣE'_j and ΣE''_j stable and non-zero against the
%        indecomposability; Lemme 2, the u_i are bijective. The head of the page
%        is cancelled throughout.
%   194  Conclusion, the u_i bijective and the E_i of dim 1; the representation
%        noted e_{n,k+1} with e_{n,1} = σ_n; the Proposition, the
%        indecomposables of G are the e_{n,k} (n ∈ Z, k ≥ 1), with
%        e_{n,k} = σ_n ⊗ e_{0,k}, e_{0,k} on k[Z]/Z^k, X the multiplication by Z
%        and Y·Z^α = αp Z^α; ě_{n,k} = e_{−n−p(k−1),k}.
%   195  the homomorphisms A(H) → A(G) → A(F); A(H) generated by the σ_n with
%        σ_n σ_m = σ_{n+m}, embedded by j*(σ_n) = e_{n,1}; A(F) with basis the
%        e_n and A(F) = Z[e_2], receiving i*(e_{n,k}) = e_k; Λ = A(H), ξ = e_{02},
%        and the aim A(G) = Λ[ξ], i.e. (1) ξ^n = Σ_1^{n+1} c_{ni} e_{0i} with
%        c_{n,n+1} = 1.
%   196  applying i* gives e_2^n = Σ d(c_{ni}) e_i against
%        e_2^n = Σ_1^{n+1} a_{ni} e_i with a_{n,n+1} = 1, so ε(c_{ni}) = 0 for
%        i > n+1 and ε(c_{n,n+1}) = 1; since the c_{ni} lie in Λ⁺,
%        c_{ni} = 0 for i > n+1 and c_{n,n+1} is a σ_α; the weight computation on
%        σ_α e_{0,p+1} = e_{α,p+1} begins.
%   197  the highest weight in ξ^n gives pα ≤ pn, hence α ≤ 0; the weights in ξ
%        being 0 and p gives α ≥ 0; so α = 0, σ_α = 1, c.q.f.d. The closing
%        Proposition: A(G) is the polynomial algebra over Λ = Z(Γ) (Γ = Z)
%        generated by ξ = e_{02}, the augmentation of Λ the usual one,
%        d(ξ) = 2, ∨ the usual involution on Λ, and ξ̌ = e_{−p,2} = σ_{−p}·ξ.

\documentclass[11pt,a4paper]{article}
% Le préambule commun à toutes les transcriptions du fonds Grothendieck.
%
% Il tient en une idée : **l'incertitude se compose, elle ne se résout pas.**
% Une transcription qui lisse un mot douteux a détruit la seule chose pour
% laquelle on la faisait — un lecteur doit pouvoir savoir, à chaque endroit,
% ce qui était sur le feuillet et ce qui a été ajouté. Les six macros qui
% suivent sont donc l'appareil critique, et non de la décoration.
%
% Les mêmes commandes sont reconnues par `scripts/render.mjs`, qui produit la
% vue de lecture du site. Une macro ajoutée ici doit l'être aussi là-bas,
% faute de quoi le rendu échouera bruyamment — ce qui est voulu : un rendu
% silencieusement faux est pire qu'un rendu absent.

\ProvidesPackage{grothendieck}[2026/08/08 Transcriptions du fonds Grothendieck]

\usepackage{amsmath,amssymb,amsthm}
\usepackage{mathrsfs}
\usepackage{stmaryrd}
% Les diagrammes commutatifs sont partout dans ce fonds : c'est la langue
% même de la géométrie algébrique de Grothendieck, et une transcription qui
% les rendrait en prose ne transcrirait rien.
\usepackage{tikz-cd}
% Français par défaut — c'est la langue des feuillets — mais l'anglais est
% chargé aussi : la traduction et le résumé sont des documents anglais, et la
% typographie française (espaces avant les deux-points) y serait une faute.
% Ces documents basculent par \selectlanguage{english} en tête de corps.
\usepackage[english,french]{babel}
\usepackage{fontspec}
\usepackage{geometry}
\geometry{a4paper,margin=25mm,marginparwidth=28mm}
\usepackage{marginnote}
\usepackage{xcolor}
% ulem, not soul. `soul`'s \st and \ul are built on a character-by-character
% scan that XeLaTeX's font machinery breaks: the document fails with an
% undefined control sequence rather than degrading. ulem does the same two jobs
% and is Unicode-engine safe. `normalem` stops it hijacking \emph, which the
% transcriptions use for Grothendieck's own emphasis.
\usepackage[normalem]{ulem}
\usepackage{hyperref}
\hypersetup{colorlinks=true,linkcolor=black,urlcolor=black,pdfborder={0 0 0}}

% --- Métadonnées du lot -----------------------------------------------------
% Reprises telles quelles par le script de rendu, qui les affiche en tête de
% la vue de lecture. Elles fixent ce que le fichier prétend couvrir : sans
% elles, une transcription flotte sans point d'attache dans un fonds de seize
% mille feuillets.

\newcommand{\folder}[1]{\gdef\grothFolder{#1}}
\newcommand{\batch}[1]{\gdef\grothBatch{#1}}
\newcommand{\leaves}[2]{\gdef\grothFirst{#1}\gdef\grothLast{#2}}
\newcommand{\foldertitle}[1]{\gdef\grothFoldertitle{#1}}
\gdef\grothFolder{?}\gdef\grothBatch{?}\gdef\grothFirst{?}\gdef\grothLast{?}\gdef\grothFoldertitle{}

% --- Le résumé --------------------------------------------------------------
%
% Il ouvre la lecture modernisée, sous le titre « Résumé », et il est composé
% en petit. Ce n'est pas un ornement : c'est le seul passage du document qui
% ne soit pas des mathématiques, et un lecteur doit pouvoir le reconnaître
% avant de l'avoir lu — puis le sauter s'il connaît déjà le sujet, ou s'y
% arrêter s'il ne le connaît pas. Le filet qui le suit marque la reprise du
% corps.

\newenvironment{resume}
  {\small\setlength{\parskip}{0.45em}}
  {\par\medskip\hrule\medskip}

% --- Mots-clés ---------------------------------------------------------------
%
% En anglais, en clôture du résumé de la lecture modernisée : le vocabulaire
% moderne sous lequel chercher ce que les feuillets construisent. Le manifeste
% du site (`npm run manifest`) les extrait pour étiqueter la cote entière —
% cette ligne est la source unique des tags, il n'y a pas de fichier à côté.
\newcommand{\keywords}[1]{\par\smallskip\noindent
  {\footnotesize\textsc{Keywords} --- \textit{#1}}\par}

% --- L'appareil critique ----------------------------------------------------

% Le numéro de feuillet, en marge. C'est l'ancre qui permet de régler un
% désaccord sur une formule en regardant *un* feuillet plutôt qu'une chemise
% de six cents. Le site s'en sert aussi pour tourner les pages du fac-similé
% quand on fait défiler la transcription.
\newcommand{\leaf}[1]{%
  \marginnote{\footnotesize\color{gray!70}\textsf{#1}}%
  \label{leaf:#1}\ignorespaces}

% Illisible : on ne devine pas. Un mot inventé qui se lit comme les autres est
% le pire résultat possible.
\newcommand{\ill}{\textcolor{red!60!black}{[\,\ldots\,]}}

% Lecture douteuse : le mot est donné, et signalé comme incertain.
\newcommand{\uncertain}[1]{\uline{#1}}

% Ajout de l'éditeur — une lettre manquante, un mot sous-entendu.
\newcommand{\add}[1]{\textcolor{blue!50!black}{[#1]}}

% Biffé par Grothendieck lui-même. Ce qu'il a rayé fait partie du document :
% une rature dit souvent où la pensée a changé de direction.
\newcommand{\struck}[1]{\sout{#1}}

% Note du transcripteur, sur l'état matériel du feuillet.
\newcommand{\note}[1]{\par\smallskip{\small\color{gray!60!black}\textit{[#1]}}\par\smallskip}

% Note marginale de Grothendieck, distincte du corps du texte.
\newcommand{\marginal}[1]{\par\smallskip\noindent\hspace{1em}%
  {\small\color{gray!60!black}#1}\par\smallskip}

% --- Titre ------------------------------------------------------------------

\AtBeginDocument{%
  \begin{center}
    {\large\bfseries Fonds Alexandre Grothendieck --- cote n\textsuperscript{o}~\grothFolder}\\[2pt]
    {\itshape\grothFoldertitle}\\[4pt]
    {\small feuillets \grothFirst--\grothLast\ (lot \grothBatch)}\\[2pt]
    {\footnotesize\color{gray!60!black}Transcription établie d'après le fac-similé numérisé
     par l'Université de Montpellier.\\
     Les lectures incertaines sont soulignées, les illisibles notées [\,\ldots\,],
     les ajouts entre crochets.}
  \end{center}
  \vspace{1em}}

\endinput


\folder{10}
\batch{10}
\pages{181}{197}
\dating{[à partir de 1958]}
\watermark{Édition de démonstration}
\foldertitle{Catégories tensorielles : notes manuscrites (s.d.), tapuscrits (s.d.)}

\begin{document}

\section*{L'hypergroupe discret des classes indécomposables (suite)}

\note{Les pages 181 et 182 continuent sans coupure la rédaction commencée page
180. Elles sont écrites d'un trait rapide : les formules et les énoncés
soulignés portent, une grande partie de la prose de liaison ne porte pas.}

\page{181}

sur $\Gamma^{+}$, i.e.\ la fonction « dimension d'un module ».

\note{Le symbole est tracé d'une plume chargée, un $\Gamma$ à pied empatté ; il
désigne le groupe libre $\mathbf{Z}(I)$ de la page 180, dont $\Gamma^{+}$ est le
cône positif. La même lettre revient page 182 et, page 180, y avait été laissée
\ill{} sous une tache d'encre.}

Dans le cas où \ill{} un $e$, on a $n(e) = 1$, et si \ill{} un $\vee$, on a
\[ n(\check{\imath}) = n(i) . \]

$\mathcal{R}$ est bien définie si, \ill{} donne de $I$ et de la fonction $n(i)$,
\ill{} se donne des représentants $M_{i}$ \ill{} de degré $i$ avec \ill{}
\uncertain{module} \ill{} par $k^{n(i)}$ et pour $i, j$ donnés, un \ill{}
\uncertain{groupe} $\operatorname{Hom}_{j,i}\bigl(k^{n(i)}, k^{n(j)}\bigr)$, de
façon que $\operatorname{Hom}_{kj} \times \operatorname{Hom}_{ji}$ soit \ill{}
par composition dans $\operatorname{Hom}_{ki}$, etc.\ldots

Supposons que dans $\mathcal{R}$ on ait une \uncertain{notion} de produit
tensoriel, alors $A(K)$ est muni d'une loi d'anneau \add{\ill{}}
\note{Un mot inséré au-dessus de « anneau » ne se lit pas ; il se termine en
« -if ».}
et commutatif, qui sur $\Gamma^{+}(K)$ y correspond à la loi de produit
tensoriel. $n$ est une fonction \uncertain{multiplicative} sur $\Gamma^{+}$, et
la compatibilité avec $e$ \struck{signifie} \add{indique} que $e$ est une unité
pour \ill{} la loi d'algèbre.
\note{Le mot que $\ill{}$ remplace est souligné sur la page.}
La compatibilité avec $\vee$ \struck{signifie} \add{indique} que $\vee$ est un
automorphisme de l'algèbre $A$. Enfin la compatibilité simultanée de $e$,
$\vee$, $\otimes$ indique que $i * \check{\imath}$ a toujours pour composante
une unité suivant $e$. $I$ est alors une \emph{« hypergroupe discret »}.
\note{L'accord est celui de la page.}

\page{182}

Si tous $M \in \mathcal{C}$ est complètement réductible, et $\mathcal{C}$ est
engendré par un \struck{\ill{}} $G$ avec une \uncertain{algèbre} $A$ sur $G$,
alors $\Gamma$ s'identifie : l'algèbre des \struck{\ill{}}
\struck{combinaisons} linéaires des \ill{} ( fonctions ) sur

\begin{itemize}
\item[a)] ses représentations, dimensions de $G$,
\end{itemize}

$\Gamma^{+}$ : \add{l'ensemble des} \struck{\ill{}} ( fonctions des
représentations linéaires de $G$ ( avec conservation de toutes les structures :
$n(f)$ devient $f(e)$ $(e \in G)$, $\vee$ devient $\vee$, $e$ devient la
fonction $1$ sur $G$, \ill{} précédemment ; noter d'ailleurs qu'il s'agit de
fonctions centrales ). Si \uncertain{p. ex.} $G$ est un groupe algébrique
semi-simple \ill{} \ill{}, on obtient une \ill{} structure concrète de $\Gamma$,
qui est l'algèbre des fonctions \ill{} sur $T/W$, $T$ un « tore » maximal, $W$
le groupe de Weyl ( il faudrait essayer de trouver une \ill{} intrinsèque sinon
des \ill{} irréductibles, i.e.\ des \ill{} $\in I$ ; pour \ill{} il faut mettre
sur $T/W$ une \ill{} de hypergroupe, dont on \ill{} \uncertain{construire} les
\ill{} irréductibles ).
\marginal{?}
\note{Le point d'interrogation est porté dans la marge de gauche, en regard de
cette parenthèse.}

Cas de $G = k^{*}$ ( $k$ alg.\ clos pour simplicité ). Alors $\Gamma(G) \simeq$
algèbre du groupe $\mathbf{Z}$ ( $\mathbf{I}(G) \simeq \mathbf{Z}$, avec \ill{},
les \ill{} ).

\section*{Le groupe additif : $A(k)$}

\note{Les pages 183 à 188 traitent le cas du groupe additif ; elles ne portent
pas de numéro de partie sur les feuilles, au contraire des deux suivantes, qu'il
numérote lui-même II (page 189) et III (page 191).}

\page{183}

\marginal{Caractéristique $0$ ?}

Cas de $G = k$ ( $k$ alg.\ clos pour simplifier ). Toutes les représentations
rationnelles de $k$ \uncertain{sont} unipotentes, leur classification est donc
identique à celle des opérateurs unipotents, ou encore des opérateurs
nilpotents, i.e.\ des modules sur $k[X]$ qui sont annulés par une puissance
convenable de $X$. Donc les \uncertain{termes} élémentaires sont de type $X^{k}$,
donc $I \simeq \mathbf{N}^{+}$ ( le \ill{} à $i$ \ill{} $\in \mathbf{N}^{+}$
étant $k[X]/X^{i}$ ). L'élément neutre est $1 \in \mathbf{N}^{+}$, l'involution
l'identité. Quelle est la multiplication ?

\page{184}

\marginal{\ill{} ???}

Soit $k$ un corps. $k$ est un \uncertain{groupe algébrique} dont \ill{} nous
cherchons les représentations linéaires, i.e.\ les \emph{applications
polynomiales} $P(x)$ à valeurs dans un \ill{}, telles que
\[ P(x+y) = P(x)\,P(y) . \]

Les \emph{représentations indécomposables}. Alors le commutant de l'ens.\ des
$P(x)$ est un corps, \struck{\ill{}} soit $D$ et les $P(x)$ prennent \ill{}
valeurs dans le \uncertain{centre} $Z$ de $D$.
\add{( comme précédemment, on \uncertain{voit} \ill{} que $D = Z$ )}

$Z$ est un corps \add{\uncertain{commutatif}} \ill{} sur $k$, donc \ill{} une
application $k$-polynomiale de $k$ dans une \struck{\ill{}} \ill{} commutative
$Z$ \uncertain{de} $k$, telle que $P(x+y) = P(x)P(y)$. \struck{Soit \ill{}}
\ill{} $P(x) = 1$ pour tout $x$ ; \struck{\ill{} le corps des}

\struck{\ill{} relatives à la clôture \uncertain{algébrique} $\hat{k}$ de $k$,
on \ill{} obtient $x \mapsto \hat{P}(x)$ hom.\ ( \ill{} ) de $\hat{k}$ dans le
groupe multiplicatif des \ill{} inversibles de $Z$, $\hat{k}$-algèbre de
dim.\ \ill{} \ill{} } \qquad $A = \hat{k} \otimes_{k} Z$

\struck{$P(x) = a_{n} x^{n} + a_{n-1} x^{n-1} + \ldots + a_{1} x + 1$},
\struck{$(a_{i} \in Z)$}, \struck{avec $a_{n} \neq 0$} ;

\struck{$P(x+y) = a_{n} x^{n} + a_{n} y^{n} + \ldots$}

\struck{$P(x)P(y) = a_{n}^{2} x^{n} y^{n} + \ldots$}

\note{Ce bloc est enfermé dans un cadre tracé à la main et barré de deux longues
diagonales ; il rend ses formules et presque rien de la prose entre elles.}

\ill{} un polynôme de degré $n > 0$, $P(x+y)$ \ill{} de degré $n$ et
$P(x)P(y)$ de degré $2n$, absurde.

\page{185}

Ainsi $P(x) = 1$, ce qui \ill{} que $Z = k$, donc

\emph{Prop.} \struck{Toute} \add{tt les} représentations \add{lin.}
rationnelles indécomposables \ill{} $k$ \ill{} de dim $1$.

\emph{Corollaire} : Toute représentation linéaire \uncertain{rationnelle} de $k$
\ill{} \ill{} \ill{} mettre sous forme \struck{diagonale} triangulaire, avec des
$1$ sur la diagonale.

Si \ill{} en caractéristique $p \neq 0$ \ill{} \struck{\ill{}} rationnelles de
$k$ \uncertain{type} $F$, $\lambda \mapsto x^{p}$, \struck{\ill{}} d'où une
\ill{} sur les classes de représentations linéaires, en comparant avec $F$ et
ses itérés ; la \uncertain{situation} est tout à fait distincte de ce qui se
passe en car.\ $0$, \ill{} les repr.\ de dim $2$ \ill{} \ill{} \struck{\ill{}}
\ill{} des \ill{}
\[ f(\lambda,\mu)(x) = \lambda x^{p^{h}} + \mu \qquad
   ( h \text{ fixé, } x \text{ l'indéterminée} ) : \]
\begin{align*}
\tau_{y} f(\lambda,\mu)(x)
 &= \lambda (x+y)^{p^{h}} + \mu \\
 &= \lambda x^{p^{h}} + \mu + \lambda y^{p^{h}}
  = f\bigl(\lambda,\ \mu + \lambda y^{p^{h}}\bigr) , \quad\text{i.e.}
\end{align*}
\[ E_{h}(y) = \begin{pmatrix} 1 & y^{p^{h}} \\ 0 & 1 \end{pmatrix} ; \]
ces représentations \ill{} sont \struck{\ill{}} \ill{} distinctes, \ill{}
\uncertain{inéquivalentes} deux à deux, d'où une \ill{} de classes \ill{} bien
\ill{} distinctes. En car.\ $0$, il n'y en a que \uncertain{deux}, \ill{}

\note{Une coulure d'encre traverse les trois dernières lignes de la page.}

\page{186}

\marginal{$k$ de car$^{0}$}

\marginal{La classification \ill{} \ill{} classes de \ill{} \ill{} d'invariants
\ill{} \ill{} dépend \emph{aussi} de la caractéristique. \ill{}}
\note{Cette note court en biais dans la marge inférieure gauche ; deux de ses
groupes de mots sont soulignés.}

Supposons maintenant la caractéristique nulle. Alors les représentations
\ill{} rationnelles des \struck{\ill{}} \emph{groupes} \emph{abéliens} \ill{}
correspondent à \struck{certaines} représentations lin.\ de l'algèbre de Lie
\emph{abélienne} $k$ ; il résulte de \ill{} qu'il s'agit exactement des
représentations \emph{nilpotentes} de l'algèbre $k$, dont la classification est
celle des $k[X]$-modules de type fini $X$-primaires. Ainsi

\textbf{Prop.\ 2} La classification des repr.\ linéaires des \ill{}
alg.\ \emph{abéliens} \ill{}, via les \ill{} \ill{}, équivalente : la
classification des $k[X]$-modules de type fini $(X)$-primaires.

\note{Un trait vertical porté dans la marge accompagne cet énoncé, en regard
duquel se lit « $k$ de car$^{0}$ ».}

Les $k[X]$-modules $(X)$-primaires \emph{indécomposables} sont ceux qui sont de
type $k[X]/X^{n}$ $(1 \leq n < +\infty)$. \ill{} voit \uncertain{aussitôt} que
la \struck{\ill{}} \ill{} de $k[X]$-modules $E_{n}$ \ill{}
\uncertain{isomorphe} à $E$ \ill{} Ainsi

\textbf{Corollaire} L'anneau \add{$A(k)$} des classes de représent\ill{} des
\ill{} alg.\ \ill{} \ill{} bien \ill{}

\page{187}

les classes $E_{n}$ des modules $k[X]/X^{n}$, $n = 1, 2, \ldots$ . L'involution
$\vee$ dans $A$ \ill{} à l'identité. On a
\[ d(E_{n}) = n , \qquad E_{1} \text{ est l'unité} . \]

Structure multiplicative. Posons $\xi = E_{2}$, classe de $k[X]/(X^{2})$. Alors
$\xi^{n}$ est la classe de
\[ F_{n} = k[X_{1}, \ldots, X_{n}]\,/\,(X_{1}^{2}, \ldots, X_{n}^{2}) , \]
muni de l'opérateur nilpotent $T$ de multiplication par la classe $\dot{S}$ de
$S = X_{1} + \ldots + X_{n}$. Les \ill{} itérés $T^{k}$ de $T$ \ill{} la
multiplication par
\[ \dot{S}^{k} = k! \sum_{1 \leq i_{1} < \ldots < i_{k} \leq n}
      \dot{X}_{i_{1}} \ldots \dot{X}_{i_{k}} . \]

\note{Une première leçon de cette formule la précède sur la même ligne, barrée :
$\dot{S}^{k} = k!\,\sum$ avec « $H \subset [1,n]$ » et « $\# H$ » portés sous le
signe de sommation, la sommation sur les parties de $[1,n]$ remplacée par la
sommation sur les suites croissantes d'indices.}

En particulier $T^{n}$ est la multiplication par
$\dot{S}^{n} = n!\, \dot{X}_{1} \ldots \dot{X}_{n}$, et $T^{n+1}$ est nul.

On en conclut aussitôt que $T^{n} F_{n}$ est de dim $1$ ( base formée de
$\dot{X}_{1} \dot{X}_{2} \ldots \dot{X}_{n}$ ) et $T^{n+1}(F_{n}) = 0$ ; d'où
\ill{}, tout \uncertain{d'abord}, que \struck{\ill{}} $\xi^{n}$ est \ill{}
combinaison linéaire des $E_{1}, \ldots, E_{n}$ à coefficients
$\in \mathbf{Z}^{+}$, celui de $E_{n}$ étant $1$ :

\page{188}

\[ \xi^{n} = L_{n}(E_{1}, \ldots, E_{n}) = \sum_{i=1}^{n} c_{ni} E_{i}
   \qquad \underline{c_{nn} = 1} \]

On en conclut la

\textbf{Proposition 3} $A(k)$ est l'anneau de polynômes $\mathbf{Z}[\xi]$ sur
$\mathbf{Z}$ engendré par $\xi = E_{2}$ $\bigl( d(\xi) = 2 \bigr)$

N.B. Le calcul explicite donne
\begin{align*}
\xi^{0} &= E_{1}                    &  E_{1} &= 1 \\
\xi^{1} &= E_{2}                    &  E_{2} &= \xi \\
\xi^{2} &= E_{1} + E_{3}            &  E_{3} &= \xi^{2} - 1 \\
\xi^{3} &= 2 E_{2} + E_{4}          &  E_{4} &= \xi^{3} - 2\xi \\
\xi^{4} &= 2 E_{1} + 3 E_{3} + E_{5} & E_{5} &= \xi^{4} - 3\xi^{2} + 5 \\
        &\qquad \ldots               &        &\qquad \ldots
\end{align*}

\note{La dernière ligne de la colonne de droite est bien « $+\,5$ » sur la
page ; la relation $\xi^{4} = 2E_{1} + 3E_{3} + E_{5}$ de la colonne de gauche,
avec $E_{1} = 1$ et $E_{3} = \xi^{2} - 1$, donnerait $+\,1$. La leçon est donnée
telle quelle.}

( Toutes les structures de \ill{} sont bien \ill{} \ill{} dans un \ill{} de
\uncertain{structures}, \emph{2} \ill{} la structure d'anneau, qui est un peu
cachée ! )
\note{Le premier mot illisible de cette parenthèse commence par une capitale
$A$ ; il ne se lit pas.}

\section*{II. \quad $A(k^{*})$}

\page{189}

\marginal{$k$ \ill{} : $(p)$}

II \qquad $A(k^{*})$. \struck{\ill{}} \quad $^{\circ}$) Représentations
irréductibles.

Se font via une représentation rationnelle de $k^{*}$ dans $L^{*}$, où $L$ est
un corps extension finie de $k$.
\[ P(\lambda) = \sum_{-n}^{n} \lambda^{i} c_{i} \qquad ( c_{i} \in L ) \]
on a $P(\lambda\lambda') = P(\lambda) P(\lambda')$, i.e.
\[ \sum_{i} \lambda^{i} \lambda'^{i} c_{i}
 = \sum_{i,j} \lambda^{i} \lambda'^{j} c_{i} c_{j} \ , \]
\note{Un symbole suit $c_{i} c_{j}$ sur la ligne et ne se lit pas.}
et comme c'est une identité en $\lambda, \lambda'$ on doit avoir
\[ c_{i} = c_{i}^{2} , \qquad c_{i} c_{j} = 0 \ \text{ si } i \neq j , \]
donc il ne peut y avoir qu'\struck{\ill{}} \ill{} des $c_{i}$ qui soit
$\neq 0$, soit $c_{n}$, et le \ill{} $c^{2}$ \ill{} exige $c_{n} = 1$. Donc
$P(\lambda) = \lambda^{n}$ ; l'irréductibilité \ill{} exige $L = k$. Ainsi

\textbf{Proposition} Les seules repr.\ \struck{irréd.} irréductibles de $k^{*}$
( $k$ corps \uncertain{quelconque} ) sont les repr.\ de dim $1$,
\uncertain{définies} \ill{} aux repr.\ \struck{\ill{}}
\[ \sigma_{n}(\lambda)\, t = \lambda^{n} t . \]
\note{Le symbole opéré, au premier membre, est tracé d'une plume bouclée et ne
se laisse pas identifier avec certitude au $t$ du second membre.}

\textbf{Corollaire}
\[ \sigma_{n} \sigma_{m} = \sigma_{n+m} , \qquad
   \check{\sigma}_{n} = \sigma_{-n} \]

Plus généralement, considérons une repr.\ linéaire \add{rationnelle}
\struck{quelconque} de $k^{*}$ dans un espace \ill{}. On \ill{} la \ill{} \ill{}
sous la forme

\page{190}

\[ P(\lambda) = \sum_{0}^{n} c_{i} \lambda^{i} \qquad ( c_{i} \in L(V) ) \]
on voit comme ci-dessus que
\[ c_{i} c_{j} = 0 \ \text{ si } i \neq j , \qquad c_{i}^{2} = c_{i} \]
De plus $\sum c_{i} = P(1) = 1$, donc les $c_{i}$ sont des projecteurs deux à
deux orthogonaux, \ill{}

\textbf{Proposition} Toute représentation \add{lin.\ rat.} de
\uncertain{$C^{*}$} est complètement réductible ( $k$ corps \ill{}
quelconque ).
\note{La lettre porte clairement une capitale $C$ ; le contexte, et le
« $A(k^{*})$ » du corollaire qui suit, appelleraient $k^{*}$. La leçon est
donnée telle quelle.}

\textbf{Corollaire} $A(k^{*}) \simeq \mathbf{Z}(\mathbf{Z})$, $\mathbf{Z}$
\ill{} \ill{} \ill{} rationnels, dans \ill{} \ill{}, $\sigma_{n}$ \ill{} \ill{}
de $\mathbf{Z}(\mathbf{Z})$, \ill{} \ill{} \ill{} cf.
\note{Le symbole entre les parenthèses de $\mathbf{Z}(\ )$ est repassé à la
plume jusqu'à la tache ; il revient trois fois dans la même phrase.}

L'involution \ill{} \ill{} compatible avec l'augmentation, \ill{}
\emph{d'anneau} \ill{} $\vee$.

\add{le \uncertain{rang}} \ill{} \ill{} \ill{} $\lambda^{n}$ : \ill{}

\emph{Généralisation aux « tores » \uncertain{isomorphes} \ill{}}

Si $\Gamma$ \ill{} le \uncertain{groupe} des \ill{} r.t.\ de $G$ \ill{} $k^{*}$,
$\Gamma$ \ill{} $\simeq \mathbf{Z}^{n}$, et $A(G) \simeq \mathbf{Z}(\Gamma)$.

\section*{III. \quad $G = k \cdot k^{*}$, produit semi-direct}

\page{191}

III \qquad $G = k \cdot k^{*}$ produit semi-direct, $k^{*}$ opérant dans $k$
par
\note{Les deux facteurs portent chacun, écrit dessous entre guillemets, le nom
qui les désigne dans la suite : « $F$ » sous $k$, « $H$ » sous $k^{*}$.}
\[ \sigma_{p}(\lambda) t = \lambda^{p} t , \]
( \uncertain{notation} \ill{} \emph{\ill{}} ). \ill{} \ill{} \ill{}
car.\ $0$ et l'algèbre de Lie de \ill{} \ill{} est extension de $F = k$ par
$H = k$, $H$ opérant sur $F$ par \ill{} ( $X$ et $Y$ \ill{} les \ill{} \ill{}
\uncertain{canoniques} de $F$ et $H$ ) :
\begin{equation*}
(1) \qquad [\,Y, X\,] = p X
\end{equation*}

Considérons une représentation linéaire indécomposable rationnelle $\varphi$ de
\struck{\ill{}} $G$, soit $\varphi'$ \ill{} \uncertain{différentielle}. \ill{}
\ill{} \ill{} \ill{} propres \ill{} \ill{} : $H \subset G$ )
\[ V = \sum V_{n} \]
( $Y . a = n\, a$ \ \ si \ $a \in V_{n}$ )
\note{Un $p$ barré précède le $n$.}. La condition $(1)$
\ill{} que les \ill{} $X_{ji} \in L(V_{i}, V_{j})$ de $X$
\note{Le second indice de $X_{ji}$ est écrit deux fois, la première barrée.}
\uncertain{vérifient} :
\[ (j - i)\, X_{ji} = p\, X_{ji} \]
d'où $X_{ji} = 0$ \uncertain{sauf} si $j - i = p$.

\page{192}

Si donc les $X_{ji}$ sont nuls, alors ( $V$ étant indécomposable ) il \ill{}
qu'un seul $V_{n}$ non nul, qui est de dimension un. La \ill{} de $V$ est ainsi
obtenue ( \ill{} \ill{} repr.\ $\sigma_{n}$ de $H = C^{*}$ \ill{} \ill{}
$G \to H$ ) par $\sigma_{n}$. S'il existe un $X_{ji}$ non nul, alors
\add{$p - a = p$, \ill{}} \ill{}
\[ V = V' + V'' , \]
i.e.
\[ V' = \sum_{\substack{n \equiv \alpha \\ (p)}} V_{n} , \qquad
   V'' = \sum_{\substack{n \not\equiv \alpha \\ (p)}} V_{n} \]
\ill{} stables \ill{} \ill{}, et $V' \neq 0$, donc $V' = V$.

\ill{} \ill{} donc, si \struck{\ill{}}
\add{$V_{n}, V_{n'} \neq 0 \,:\, n' - n \equiv 0 \ (p)$}
\add{Soit $n$ le plus} \struck{\ill{} \ill{} des} \ill{} $V_{n} \neq 0$, donc
les $V_{m}$ \ill{} \ill{} \ill{} \ill{} de la forme
\[ V_{n+pk} \qquad k = 0, 1, \ldots \ . \]
Soit $k$ le plus grand \ill{} tel que $V_{n+pk} \neq 0$, \struck{\ill{}} \ill{}
$k \geq 1$. Soit
\[ u_{i} = X_{n+p(i+1),\, n+pi} , \qquad E_{i} = V_{n+pi} , \]
\note{Une ligne barrée s'interpose entre les deux, dont seul un $V$ se lit.}
\ill{} \ill{} \ill{} d'applications
\[ E_{0} \xrightarrow{\ u_{0}\ } E_{1} \xrightarrow{\ u_{1}\ } E_{2}
   \ \ldots \ \xrightarrow{\ u_{k-1}\ } E_{k} . \]

\textbf{Lemme 1.} Les $u_{i}$ sont \emph{\uncertain{injectifs}}.
\add{$(0 \leq i \leq k-1)$}
( \ill{} \ill{} le \uncertain{prouver} par \uncertain{récurrence} sur $i$,
\ill{} trivial pour $u_{-1} : 0 = E_{-1} \to E_{0}$. Supposons

\page{193}

\add{$(0 \leq \ell \leq k-1)$}
\ill{} \ill{} pour $u_{\ell-1}$ \struck{\ill{}} \ill{} le \ill{} $u_{\ell}$.

\struck{\ill{}} Soit $x$ un élément du $\operatorname{Ker} u_{\ell}$,
\struck{\ill{}}, \struck{\ill{} \ill{} \ill{} \ill{}} \add{\ill{}}
$E_{\ell'}$ $(\ell' \leq \ell)$

\struck{\ill{} : \ill{} \ill{} de $E_{\ell}$ \ill{} $+\, 0$}
$(-1 \leq \ell' < \ell)$

\struck{\ill{} \ill{} \ill{} \ill{} \ill{}} \struck{\ill{}} $\ell'$
\struck{$(\ill{} \leq \ell)$}

\ill{} \ill{} il \uncertain{existe} un \uncertain{premier} \ill{} \ill{}
$x \notin E_{\ell'}$ $-\, E_{\ell'}$ \ill{} \ill{} \ill{} \ill{} de
$E_{\ell}$.

\note{Le haut de la page — six lignes — est raturé de part en part, par une
quinzaine de traits et deux longues barres horizontales ; il rend quelques
symboles et rien de la prose entre eux.}

Soit $D$ le \ill{} \uncertain{engendré} par $x$, \struck{\ill{}}
$(D \subset E_{\ell+1})$ et soit
\marginal{$E_{\ell'}\ E_{\ell'+1} \ldots E_{\ell}$}
\struck{$E_{j}$} $F$ un \uncertain{supplémentaire} de $E_{\ell}$ \ill{} \ill{}
$E_{\ell'}$ \ill{} \ill{} de $D$ ; soit pour $0 \leq j \leq k$,

\note{Une première leçon barrée précède celle qui suit, portée en accolade :
$E'_{j} = E_{j}$ si $j \leq \ell'$, $E'_{j} = D$ si $j > \ell$.}

\[ E'_{j} = E_{j} \ \text{ si } \ j \leq \ell' , \qquad
   E'_{j} = E_{j} \cap F \ \text{ si } \ \ell' \leq j \leq \ell \]

\note{Entre les deux, une troisième leçon barrée, « $E'_{j} = \ill{}$ si
$j > \ell$ ».}
\[ E''_{j} = 0 \ \text{ si } \ j \leq \ell' \ \text{ ou } \ j > \ell ,
   \qquad E''_{j} = D \ \text{ si } \ \ell' \leq j \leq \ell \]

\ill{} $E_{j} = E'_{j} \oplus E''_{j}$ pour \ill{} $j$, \ill{}
\struck{\ill{}} \add{et $\sum E'_{j}$ \ill{}} $\sum E''_{j}$ \ill{} stables
\ill{} \struck{\ill{}} \ill{} $\neq 0$, donc $\sum E_{i}$ \ill{} \ill{} \ill{}
\uncertain{indécomposable}.

\textbf{Lemme 2} Les $u_{i}$ $(0 \leq i \leq k-1)$ sont \emph{bijectifs}.

\ill{} \ill{} \ill{} \ill{} \ill{}, \ill{} \ill{} \ill{} \ill{} \ill{} \ill{}
\ill{} \ill{} \ill{}

\page{194}

\marginal{\ill{} $V'$ \ill{} $F$ \ill{} vect.\ \ill{}
\uncertain{indécomposable} \ill{} $e_{k+1}$ \ill{}}
\note{Cette note, de quatre lignes rassemblées par une accolade, court en biais
dans la marge de gauche ; elle ne se lit pas au delà.}

Conclusion : les $u_{i}$ \add{$(0 \leq i \leq k-1)$} \struck{\ill{}} sont des
bijections, les $E_{i}$ sont de dim.\ $1$.

La représentation de \ill{} \ill{} de $G$ \ill{} \ill{} de \ill{} est notée
$e_{n,k+1}$ $(k+1 \geq 2)$ ( \ill{} \ill{} \ill{} \ill{} \ill{}, \ill{} \ill{}
\ill{} $k = 0$, \ill{} \ill{} naturel de poser $e_{n,1} = \sigma_{n}$.

\struck{Ainsi}

\textbf{Proposition} Les représentations indécomposables de $G$ \ill{}
\struck{\ill{}} \add{( à équivalence près )} représentations
$e_{n,k}$ \struck{\ill{}}
\[ ( n \in \mathbf{Z} , \ k \in \mathbf{Z} , \ k \geq 1 ) \]
\ill{} \ill{} de \uncertain{construire} : \struck{\ill{}} $e_{n,1}$ est la
représentation $\sigma_{n}$ de dim $1$ \uncertain{induite} par la
représentation de \ill{} \ill{} de $H = C^{*}$,
\[ e_{n,k} = \sigma_{n} \otimes e_{0,k} , \]
\ill{} $e_{0,k}$ \ill{} la représentation \ill{} \ill{} \ill{}
\[ k[Z]/Z^{k} \]
\uncertain{définie} \ill{} : $X$ \ill{} l'\uncertain{opérateur} de
multiplication par $Z$ \ill{} $(Z^{k})$, et
\[ Y . Z^{\alpha} = \alpha p\, Z^{\alpha} \qquad (Z^{k}) \]
\add{\ill{} \ill{} \ill{} \ill{} de $e_{n,k}$ \ill{} $n$ et $n + p(k-1)$,
\ill{}}
( Les représentations $e_{n,k} = \sigma_{n} \otimes e_{0,k}$ sont deux à deux
inéquivalentes. $e_{0,k}$ est de \ill{} \ill{} $k$ ;
\[ \check{e}_{n,k} = e_{-n-p(k-1),\, k} . \]

\page{195}

Nous \ill{} déterminer la structure multiplicative de $A(G)$. Les \ill{}
$F \xrightarrow{\ i\ } G \xrightarrow{\ j\ } H$ \ill{} des \uncertain{homomorphismes}
\[ A(H) \xrightarrow{\ j^{*}\ } A(G) \xrightarrow{\ i^{*}\ } A(F) \]
$A(H)$ est l'algèbre \struck{\ill{}} \uncertain{engendrée} \ill{} les
$\sigma_{n}$ $(n \in \mathbf{Z})$ \ill{} \ill{} multiplication,
\[ \sigma_{n} \sigma_{m} = \sigma_{n+m} , \qquad \sigma_{0} = 1 . \]
il est \uncertain{plongé} dans $A(G)$ par $j^{*}$, \ill{}
\[ j^{*}(\sigma_{n}) = e_{n,1} \]
$A(F)$ est l'algèbre \struck{\ill{}} \ill{} \ill{} une base formée \ill{}
\ill{} $e_{n}$ $(n \geq 1)$ \ill{} $e_{1} = 1$. D'ailleurs
$A(F) = \mathbf{Z}[e_{2}]$. \ill{}
\[ i^{*}(e_{n,k}) = e_{k} \]

\uncertain{Identifions} $\Lambda = A(H)$ \ill{} une \ill{} de $A(G)$. Posons
$\xi = e_{02}$. \ill{} \ill{} de \uncertain{prouver} \ill{}
\[ A(G) = \Lambda[\xi] \]
\note{Une parenthèse de deux lignes accompagne la formule à droite ; elle
s'achève sur « de $\Lambda$ » et ne se lit pas au delà.}
i.e.\ \ill{} les $\xi^{n}$ $( n = 0, 1, \ldots )$ forment une \emph{base} de
$A(G)$ sur $\Lambda$. \ill{} les $e_{01}, e_{02}, e_{03} \ldots$ forment une
telle base, \ill{} il \uncertain{suffit} d'établir \ill{} \ill{} formules
\begin{equation*}
(1) \qquad \xi^{n} = \sum_{i=1}^{n+1} c_{ni}\, e_{0i} \qquad
   ( c_{ni} \in \Lambda , \ c_{n,n+1} = 1 )
\end{equation*}

\page{196}

On a \uncertain{a priori} \ill{}
\[ \xi^{n} = \sum_{i} c_{ni}\, e_{0i} \]
\ill{} $c_{ni} \in \Lambda^{+}$, d'où \ill{} \ill{} \ill{} par $i^{*}$ :
\[ e_{2}^{\,n} = \sum_{i} i^{*}(c_{ni})\, e_{i} \]

\struck{\ill{} \ill{} $i^{*}(c_{p}$}

Or \ill{}
\[ i^{*}(\lambda) = d(\lambda) \qquad \lambda \in \Lambda \]
\ill{} \ill{} l'augmentation de $\Lambda$, d'où
\[ e_{2}^{\,n} = \sum_{i} d(c_{ni})\, e_{i} \]
\ill{} \ill{} \uncertain{sait} d'ailleurs \ill{} \ill{}
$e_{2}^{\,n} = \sum_{1}^{n+1} a_{ni}\, e_{i}$, \ill{} $a_{ni}$ \ill{} \ill{}
\ill{} $\geq 0$ et \ill{} $a_{n,n+1} = 1$. \ill{} \ill{}
\[ \varepsilon(c_{ni}) = 0 \ \text{ si } \ i > n+1 , \qquad
   \varepsilon(c_{n,n+1}) = 1 \]
d'où, \ill{} les $c_{ni} \in \Lambda^{+}$,
\[ c_{ni} = 0 \ \text{ si } \ i > n+1 , \qquad
   c_{n,n+1} \ \text{est de la forme} \ \sigma_{\alpha} \]
( $\alpha$ \uncertain{dépendant} de \ill{} ). D'où la \ill{} $(1)$, mais \ill{}
doit \ill{} \ill{}
\[ c_{p,p+1} = \sigma_{\alpha} \]
\note{Deux leçons de la même formule précèdent celle-ci sur la ligne, toutes
deux barrées, la première « $c_{p,p+1} = \sigma$ » avec un indice illisible.}
On \ill{} \ill{} plus \uncertain{haut} \ill{} \ill{} le \uncertain{cas}.
\[ \sigma_{\alpha}\, e_{0,p+1} = e_{\alpha,p+1} \qquad \alpha + pn , \]
\ill{} \ill{} plus

\page{197}

\ill{} \ill{} \ill{} \uncertain{figurer} dans $\xi^{n}$ \ill{} \ill{} \ill{}
fois le plus \ill{} \uncertain{poids} \uncertain{figurant} dans $\xi$, \ill{}
\ill{} \ill{} \ill{} \ill{} donc
\[ p\alpha \leq pn , \qquad \text{d'où} \quad \alpha \leq 0 . \]
D'autre part, les \uncertain{poids} \uncertain{intervenant} dans $\xi$ \ill{}
$0$ et $p$, donc \ill{}, donc \ill{} dans $\xi^{n}$ \ill{} \ill{} ; il faut
donc \ill{} \ill{} \ill{} \ill{} soit $\alpha = $ le \uncertain{poids}
\uncertain{intervenant} dans $\sigma_{\alpha} e_{0,p+1}$, d'où
$\alpha \geq 0$. D'où $\alpha = 0$, i.e.\ $\sigma_{\alpha} = 1$, c.q.f.d.

Ainsi :

\textbf{Proposition} $A(G)$ est l'algèbre de polynômes sur
$\Lambda = \mathbf{Z}(\Gamma)$ $(\Gamma = \mathbf{Z})$ engendrée par
$\xi = e_{02}$ ; l'augmentation de $\Lambda$ \ill{} \ill{} augmentation
usuelle, et $d(\xi) = 2$.
\note{Un $p$ barré précède le $2$.} Sur $\Lambda$, l'involution $\vee$ \ill{}
l'involution usuelle, et \ill{} .
\[ \check{\xi} = e_{-p,2} = \sigma_{-p} \cdot \xi . \]

\end{document}
